%==============================================================
% PARTE 2: NIVEL PRINCIPIANTE
%==============================================================
\chapter{Parte 2: Nivel Principiante}

%--------------------------------------------------------------
\section{Crear un Paquete}
%--------------------------------------------------------------

\subsection*{Pregunta 9}
\textit{¿Qué comando me indica todas las dependencias que puedo usar en mis paquetes?}

\textbf{Respuesta:}\\
% Escribe tu respuesta aquí.

\begin{lstlisting}[style=bash, caption={Comando para listar dependencias disponibles.}]
% Escribe aquí el comando
\end{lstlisting}

%--------------------------------------------------------------
\section{Crear un Mensaje}
%--------------------------------------------------------------

\subsection*{Pregunta 10}
\textit{¿Qué comando se debe utilizar para saber de qué tipo es el mensaje creado?}

\textbf{Respuesta:}\\
% Escribe tu respuesta aquí.

\begin{lstlisting}[style=bash, caption={Comando para verificar el tipo de mensaje MiMensaje.}]
% Escribe aquí el comando
\end{lstlisting}

% \begin{figure}[H]
%   \centering
%   \includegraphics[width=0.85\textwidth]{parte2/imagenes/p2_msg_type.png}
%   \caption{Verificación del tipo de mensaje \texttt{MiMensaje}.}
%   \label{fig:p2_msg_type}
% \end{figure}

\subsection*{Pregunta 11}
\textit{Muestra el gráfico de entorno ROS generado usando \texttt{rqt\_graph}.}

\begin{figure}[H]
  \centering
  % \includegraphics[width=0.85\textwidth]{parte2/imagenes/p2_rqt_graph.png}
  \fbox{\parbox{0.8\textwidth}{\centering\vspace{2cm} [Captura rqt\_graph Parte 2] \vspace{2cm}}}
  \caption{Gráfico \texttt{rqt\_graph} con los nodos \texttt{nodo\_envia} y \texttt{nodo\_recibe}.}
  \label{fig:p2_rqt_graph}
\end{figure}

\subsection*{Pregunta 12}
\textit{Intenta modificar la constante nombre del mensaje, modificando el código de \texttt{nodo\_envia}. ¿Qué sucede?}

\textbf{Respuesta:}\\
% Escribe tu respuesta aquí. Explica el error o comportamiento observado.

% \begin{figure}[H]
%   \centering
%   \includegraphics[width=0.85\textwidth]{parte2/imagenes/p2_const_error.png}
%   \caption{Error al intentar modificar la constante \texttt{NOMBRE}.}
%   \label{fig:p2_const_error}
% \end{figure}

\subsection*{Pregunta 13}
\textit{¿De qué tipo es el topic que aparece?}

\textbf{Respuesta:}\\
% Escribe tu respuesta aquí. Indica el tipo del topic /mitopic.

% \begin{figure}[H]
%   \centering
%   \includegraphics[width=0.85\textwidth]{parte2/imagenes/p2_topic_type.png}
%   \caption{Tipo del topic \texttt{/mitopic}.}
%   \label{fig:p2_topic_type}
% \end{figure}

%--------------------------------------------------------------
\section{Ejercicios}
%--------------------------------------------------------------

\subsection{Ejercicio 1: Paquete \texttt{p2pkg} y mensaje \texttt{P2pkg\_mensaje}}

\textit{Crea el paquete \texttt{p2pkg} y define en \texttt{interfaz} el mensaje \texttt{P2pkg\_mensaje.msg}.}

\textbf{Descripción:}\\
% Describe los pasos realizados.

\begin{lstlisting}[style=bash, caption={Creación del paquete p2pkg.}]
cd /workspace/ros2_ws/src
ros2 pkg create --build-type ament_python --license Apache-2.0 p2pkg --dependencies rclpy std_msgs
\end{lstlisting}

\textbf{Campos del mensaje \texttt{P2pkg\_mensaje.msg}:}
\begin{lstlisting}[caption={Definición del mensaje P2pkg\_mensaje.msg.}]
int32 numero
geometry_msgs/Pose posicion
string fecha
\end{lstlisting}

\textbf{Resultado del comando de verificación:}

% \begin{figure}[H]
%   \centering
%   \includegraphics[width=0.85\textwidth]{parte2/imagenes/p2_ej1_msg.png}
%   \caption{Campos del mensaje \texttt{P2pkg\_mensaje}.}
%   \label{fig:p2_ej1_msg}
% \end{figure}

\subsection{Ejercicio 2: Nodos publisher y subscriber en \texttt{p2pkg}}

\textit{Nodo publisher \texttt{nodopub\_ejercicio2} y subscriber \texttt{nodosub\_ejercicio2} para el topic \texttt{/topic\_ejercicio2}.}

\textbf{Descripción de la implementación:}\\
% Explica la lógica del publisher y subscriber.

% \begin{figure}[H]
%   \centering
%   \begin{subfigure}[b]{0.48\textwidth}
%     \includegraphics[width=\textwidth]{parte2/imagenes/p2_ej2_pub.png}
%     \caption{Terminal del publisher.}
%   \end{subfigure}
%   \hfill
%   \begin{subfigure}[b]{0.48\textwidth}
%     \includegraphics[width=\textwidth]{parte2/imagenes/p2_ej2_sub.png}
%     \caption{Terminal del subscriber.}
%   \end{subfigure}
%   \caption{Ejecución de los nodos del ejercicio 2.}
%   \label{fig:p2_ej2_ejecucion}
% \end{figure}

% \begin{figure}[H]
%   \centering
%   \includegraphics[width=0.85\textwidth]{parte2/imagenes/p2_ej2_rqt.png}
%   \caption{Gráfico \texttt{rqt\_graph} del ejercicio 2.}
%   \label{fig:p2_ej2_rqt}
% \end{figure}
